%%%%%%%%%%%%%%%%%%%%%%%%%%%%%%%%%%%%%%%%%%%%%%%%%%%%%%%%%%%%%%%%%%%%%
%
%  GUIA DE RECURSOS — LATEX PARA RELATÓRIOS (ISEL)
%  =================================================
%  Versão  : 1.0
%
%  INSTRUÇÕES RÁPIDAS
%  ------------------
%  1. Compila este ficheiro diretamente no Overleaf com pdfLaTeX.
%  2. Partilha o PDF gerado com quem precisar de ajuda em LaTeX.
%  3. Todos os links no PDF são clicáveis.
%
%  ESTRUTURA DO FICHEIRO
%  ---------------------
%  [0] Dados do documento
%  [A] Pacotes e configurações globais
%  [B] Estilos das caixas de recursos
%  [C] Corpo do guia
%      1. Editores Online
%      2. Documentação e Referência
%      3. Comunidade e Ajuda
%      4. Templates e Código no GitHub
%      5. Ferramentas Úteis
%      6. Referência Rápida de Comandos
%
%  RECURSOS INCLUÍDOS
%  ------------------
%  → Overleaf (editor online)
%  → Overleaf Learn (tutoriais)
%  → LaTeX Wikibook
%  → CTAN (pacotes)
%  → Detexify (símbolos)
%  → TeX Stack Exchange (fórum)
%  → GitHub (templates)
%  → Tables Generator
%  → Mathcha (equações)
%  → doi2bib (referências)
%
%%%%%%%%%%%%%%%%%%%%%%%%%%%%%%%%%%%%%%%%%%%%%%%%%%%%%%%%%%%%%%%%%%%%%


% ===================================================================
% [A]  PACOTES E CONFIGURAÇÕES GLOBAIS
% ===================================================================
\documentclass[12pt]{article}

% --- Idioma e codificação ---
\usepackage[portuguese]{babel}
\usepackage[utf8]{inputenc}
\usepackage[T1]{fontenc}

% --- Página ---
\usepackage{geometry}
\geometry{
    a4paper,
    top=25mm,
    bottom=25mm,
    left=25mm,
    right=25mm
}

% --- Tipografia ---
\usepackage{setspace}
\onehalfspacing
\usepackage{microtype}

% --- Cores ---
\usepackage{xcolor}
\definecolor{linkblue}{HTML}{2563EB}
\definecolor{boxbg}{HTML}{F3F4F6}
\definecolor{boxborder}{HTML}{D1D5DB}
\definecolor{titlebg}{HTML}{1E3A5F}

% --- Hiperligações ---
\usepackage{hyperref}
\hypersetup{
    colorlinks = true,
    urlcolor   = linkblue,
    linkcolor  = black,
    citecolor  = black
}

% --- Tabelas ---
\usepackage{booktabs}
\usepackage{array}

% --- Caixas ---
\usepackage{mdframed}

% --- Numeração de secções com ponto ---
\usepackage{titlesec}
\titlelabel{\thetitle.\quad}


% ===================================================================
% [B]  ESTILOS DAS CAIXAS DE RECURSOS
% ===================================================================

% Caixa cinzenta para cada recurso
\mdfdefinestyle{resourcebox}{
    backgroundcolor  = boxbg,
    linecolor        = boxborder,
    linewidth        = 1pt,
    roundcorner      = 6pt,
    innertopmargin   = 10pt,
    innerbottommargin= 10pt,
    innerleftmargin  = 12pt,
    innerrightmargin = 12pt,
    skipabove        = 6pt,
    skipbelow        = 6pt
}

% Caixa azul escuro para o cabeçalho da página
\mdfdefinestyle{headerbox}{
    backgroundcolor  = titlebg,
    linewidth        = 0pt,
    roundcorner      = 8pt,
    innertopmargin   = 16pt,
    innerbottommargin= 16pt,
    innerleftmargin  = 20pt,
    innerrightmargin = 20pt,
    skipabove        = 0pt,
    skipbelow        = 16pt
}

% Comando auxiliar: \recurso{Título}{URL}{Descrição}
\newcommand{\recurso}[3]{%
    \begin{mdframed}[style=resourcebox]
        \textbf{#1}\\[3pt]
        \href{#2}{\textcolor{linkblue}{\small\texttt{#2}}}\\[5pt]
        \small #3
    \end{mdframed}
}


% ===================================================================
% [C]  INÍCIO DO DOCUMENTO
% ===================================================================
\begin{document}
\thispagestyle{empty}

% --- Cabeçalho da página ---
\begin{mdframed}[style=headerbox]
    \color{white}
    {\LARGE \textbf{Recursos para Relatórios em \LaTeX}}\\[4pt]
    {\normalsize Guia de referência rápida — partilha com quem precisar de ajuda}
\end{mdframed}


% -------------------------------------------------------------------
\section{Editores Online — começa aqui}
% -------------------------------------------------------------------

Não precisas de instalar nada. Usa um editor online para começar de imediato.

\recurso{Overleaf — Editor LaTeX Online (recomendado)}
    {https://www.overleaf.com}
    {O editor mais popular para LaTeX. Compila diretamente no browser,
     tem colaboração em tempo real e milhares de templates prontos a usar.
     Cria uma conta gratuita e abre o template deste relatório.}

\recurso{Overleaf — Galeria de Templates}
    {https://www.overleaf.com/latex/templates}
    {Centenas de templates prontos: relatórios, teses, apresentações,
     artigos científicos, CVs, posters, e muito mais.}


% -------------------------------------------------------------------
\section{Documentação e Referência}
% -------------------------------------------------------------------

\recurso{Overleaf Learn — Guias e Tutoriais (recomendado para iniciantes)}
    {https://www.overleaf.com/learn}
    {A melhor documentação para quem está a aprender. Cobre desde o básico
     (parágrafos, listas, figuras) até tópicos avançados (TikZ, BibTeX,
     matemática). Em inglês mas muito clara e com exemplos práticos.}

\recurso{LaTeX Wikibook — Referência Completa}
    {https://en.wikibooks.org/wiki/LaTeX}
    {Referência completa e gratuita. Útil para procurar comandos
     específicos, pacotes e exemplos detalhados de qualquer funcionalidade.}

\recurso{CTAN — Documentação Oficial de Pacotes}
    {https://ctan.org}
    {Repositório oficial de todos os pacotes LaTeX. Procura aqui a
     documentação de qualquer pacote (ex: fancyhdr, tikz, minted, amsmath).}

\recurso{Detexify — Encontra o Símbolo Matemático}
    {https://detexify.kirelabs.org}
    {Desenha o símbolo que procuras e o site mostra o comando LaTeX
     correspondente. Muito útil para equações e símbolos especiais.}


% -------------------------------------------------------------------
\section{Comunidade e Ajuda}
% -------------------------------------------------------------------

\recurso{TeX Stack Exchange — Fórum de Perguntas e Respostas}
    {https://tex.stackexchange.com}
    {O melhor sítio para tirar dúvidas. Quase tudo o que precisas já foi
     perguntado e respondido. Pesquisa antes de colocar uma nova questão.}

\recurso{Reddit r/LaTeX}
    {https://www.reddit.com/r/LaTeX}
    {Comunidade ativa onde partilham exemplos, pedem ajuda e mostram
     resultados. Boa fonte de inspiração para formatação e design.}


% -------------------------------------------------------------------
\section{Templates e Código no GitHub}
% -------------------------------------------------------------------

\recurso{LaTeX Templates — Coleção de Templates Gratuitos}
    {https://www.latextemplates.com}
    {Grande coleção organizada por categoria: relatórios, teses, cartas,
     apresentações, posters científicos, e muito mais.}

\recurso{GitHub — Templates LaTeX (pesquisa)}
    {https://github.com/topics/latex-template}
    {Repositórios públicos com templates LaTeX. Podes fazer download ou
     fork de qualquer um. Pesquisa também por "latex thesis", "latex
     report" ou "latex ISEL".}

\recurso{Awesome LaTeX — Lista Curada de Recursos}
    {https://github.com/egeerardyn/awesome-LaTeX}
    {Lista mantida pela comunidade com os melhores editores, pacotes,
     ferramentas e recursos do ecossistema LaTeX.}


% -------------------------------------------------------------------
\section{Ferramentas Úteis}
% -------------------------------------------------------------------

\recurso{Tables Generator — Cria Tabelas Visualmente}
    {https://www.tablesgenerator.com}
    {Cria tabelas como no Excel e gera o código LaTeX automaticamente.
     Suporta vários estilos incluindo booktabs (o mais profissional).}

\recurso{Mathcha — Editor Visual de Equações}
    {https://www.mathcha.io}
    {Editor visual para equações matemáticas complexas com exportação
     direta para código LaTeX. Alternativa ao Detexify para fórmulas longas.}

\recurso{doi2bib — Gera Referências BibTeX Automaticamente}
    {https://www.doi2bib.org}
    {Cola um DOI ou ID arXiv e gera a entrada BibTeX correta para
     a tua lista de referências. Poupa imenso tempo.}

\recurso{PDF24 / iLovePDF — Converter e Editar Imagens}
    {https://www.pdf24.org}
    {Converte figuras, screenshots e diagramas para formatos compatíveis
     com LaTeX (PNG, JPG, PDF).}

\newpage

% -------------------------------------------------------------------
\section{Referência Rápida — Comandos Essenciais}
% -------------------------------------------------------------------

\begin{mdframed}[style=resourcebox]
\textbf{Formatação de texto}\\[6pt]
\begin{tabular}{@{}ll@{}}
\texttt{\textbackslash textbf\{texto\}}    & \textbf{negrito}              \\
\texttt{\textbackslash textit\{texto\}}    & \textit{itálico}              \\
\texttt{\textbackslash texttt\{texto\}}    & \texttt{monospace / código}   \\
\texttt{\textbackslash underline\{texto\}} & \underline{sublinhado}        \\
\end{tabular}
\end{mdframed}

\begin{mdframed}[style=resourcebox]
\textbf{Estrutura do documento}\\[6pt]
\begin{tabular}{@{}ll@{}}
\texttt{\textbackslash section\{Título\}}       & Secção numerada           \\
\texttt{\textbackslash subsection\{Título\}}    & Subsecção numerada        \\
\texttt{\textbackslash subsubsection\{Título\}} & Subsubsecção numerada     \\
\texttt{\textbackslash label\{fig:x\}}          & Define uma etiqueta       \\
\texttt{\textbackslash ref\{fig:x\}}            & Referencia uma etiqueta   \\
\end{tabular}
\end{mdframed}

\begin{mdframed}[style=resourcebox]
\textbf{Espaçamento e páginas}\\[6pt]
\begin{tabular}{@{}ll@{}}
\texttt{\textbackslash newpage}         & Força nova página               \\
\texttt{\textbackslash clearpage}       & Nova página + flush de figuras  \\
\texttt{\textbackslash vspace\{1cm\}}   & Espaço vertical                 \\
\texttt{\textbackslash hspace\{1cm\}}   & Espaço horizontal               \\
\texttt{\% comentário}                  & Linha ignorada pelo compilador  \\
\end{tabular}
\end{mdframed}

\begin{mdframed}[style=resourcebox]
\textbf{Atalhos no Overleaf}\\[6pt]
\begin{tabular}{@{}ll@{}}
\texttt{Ctrl + Enter} & Compilar                       \\
\texttt{Ctrl + /}     & Comentar / descomentar linha   \\
\texttt{Ctrl + Z}     & Desfazer                       \\
\texttt{Ctrl + F}     & Pesquisar no código            \\
\texttt{Ctrl + G}     & Ir para linha                  \\
\end{tabular}
\end{mdframed}

\vspace{0.8cm}
\begin{center}
    \small\textcolor{gray}{
        Dúvidas? Pesquisa primeiro em\
        \href{https://tex.stackexchange.com}{\textcolor{linkblue}{tex.stackexchange.com}}
        — é muito provável que a tua pergunta já tenha resposta.
    }
\end{center}

\end{document}


%%%%%%%%%%%%%%%%%%%%%%%%%%%%%%%%%%%%%%%%%%%%%%%%%%%%%%%%%%%%%%%%%%%%%
%  SNIPPETS — copia e cola onde precisares
%%%%%%%%%%%%%%%%%%%%%%%%%%%%%%%%%%%%%%%%%%%%%%%%%%%%%%%%%%%%%%%%%%%%%
%
% -- FIGURA --------------------------------------------------------
%\begin{figure}[h]
%    \centering
%    \includegraphics[width=0.8\textwidth]{images/nome.png}
%    \caption{Descrição da figura.}
%    \label{fig:nome}
%\end{figure}
%
% -- TABELA --------------------------------------------------------
%\begin{table}[h]
%    \centering
%    \begin{tabular}{|l|c|r|}
%        \hline
%        \textbf{Coluna A} & \textbf{Coluna B} & \textbf{Coluna C} \\ \hline
%        Valor 1           & Valor 2           & Valor 3           \\ \hline
%    \end{tabular}
%    \caption{Descrição da tabela.}
%    \label{tab:nome}
%\end{table}
%
% -- LISTA COM PONTOS ----------------------------------------------
%\begin{itemize}
%    \item Primeiro item.
%    \item Segundo item.
%\end{itemize}
%
% -- LISTA NUMERADA ------------------------------------------------
%\begin{enumerate}
%    \item Primeiro item.
%    \item Segundo item.
%\end{enumerate}
%
% -- EQUAÇÃO -------------------------------------------------------
%\begin{equation}
%    f(x) = ax^2 + bx + c
%    \label{eq:nome}
%\end{equation}
%
% -- REFERÊNCIA A FIGURA / TABELA / EQUAÇÃO ------------------------
%    Ver Figura~\ref{fig:nome}
%    Ver Tabela~\ref{tab:nome}
%    Ver Equação~\eqref{eq:nome}
%
%%%%%%%%%%%%%%%%%%%%%%%%%%%%%%%%%%%%%%%%%%%%%%%%%%%%%%%%%%%%%%%%%%%%%